\begin{description}[leftmargin=2em]
\item[(B) 型。]\mbox{}\par 设 $\gamma:[0,1]\to D$ 为从 $\gamma(0)=*$ 出发的消失路径。要求 $\gamma'(0)\in\CC\setminus\R^+$，且 $\gamma|(0,1]$ 与基路径 $\chi$ 不交；从现在起，每当考虑 $\pi$ 的消失路径时，都作这些固定假设。设 $\widetilde\Delta_\gamma$ 为相应 Lefschetz 顶针 $\Delta_\gamma$ 在 $\widetilde E$ 中的原像。为其赋予任意分次 $\widetilde\alpha^\#$，以及同构意义下唯一的 Pin 结构 $\widetilde P^\#$。

为使该膜成为等变膜，还须选择一个同构 $f:\widetilde P^\#\to\iota_{\widetilde E}^*(\widetilde P^\#\otimes\lambda^{\mathrm{top}}(T\widetilde\Delta_\gamma))$。由\stref{引理}{14.2}，这类同构典范地对应于不动点集的余定向；在此情形中，该不动点集为 $V_\gamma\subset\widetilde\Delta_\gamma$。由于 $\widetilde\Delta_\gamma$ 纤维化在 $\gamma$ 的原像 $\widetilde\gamma$ 上，有
\begin{equation}\label{eq:18.21}
D\widetilde\pi(T\widetilde\Delta_\gamma|V_\gamma)=\R\widetilde\gamma'(0)=\R\sqrt{\gamma'(0)},
\end{equation}
因此，余定向自然对应于平方根的选择。我们的约定是，始终取\emph{虚部为正}的平方根。\eref{18.21}也表明，$\widetilde\Delta_\gamma$ 满足\stref{假设}{14.4}。最后，定义 $\F(\widetilde E)^{\mathrm{naive}}$ 要求为每个 $\widetilde\Delta_\gamma$ 选择实值法向角，这里取为 $\arg(\gamma'(0))/2\in(0,\pi)$。

如\secref{14d}所指出的，$\widetilde\Delta_\gamma$ 上的等变膜结构诱导其不动点集 $V_\gamma$ 上的膜结构。具体地，分次为
\[
\alpha^\#_{V_\gamma}(x)=\alpha^\#_{\widetilde\Delta_\gamma}(0,x)-\frac1{2\pi}\arg(\gamma'(0)),
\]
而 Pin 结构是在同构意义下唯一能作为 $\Delta_\gamma$ 上某个 Pin 结构的边界的结构。这意味着，当 $n=1$ 时，消失圈带有非平凡 Pin 结构，而不是平凡 Pin 结构；理由是，对任意紧连通曲面 $S$，$H^2(S,\partial S;\Z/2)\cong\Z/2$ 中的相对第二 Stiefel–Whitney 类等于模 $2$ Euler 示性数。对 $S=D$，该类非零，因此不可能存在沿边界平凡的 Pin 结构。
\end{description}
\SourcePage{287}
\StatementNumber{9}\begin{lemma}\label{st:18.9}
(i) 对任意 (B) 型等变膜 $\widetilde\Delta^\#=\widetilde\Delta_\gamma^\#$，有 $HF^*(\widetilde\Delta^\#,\widetilde\Delta^\#)\cong H^*(S^{n+1};\K)$。$\Z/2$ 作用在 $HF^0$ 上为 $\Id$，在 $HF^{n+1}$ 上为 $-\Id$。

(ii) 取消失路径对 $\gamma_k$，$k=0,1$，并设 $\widetilde\Delta_k^\#$ 为相应的 (B) 型等变膜。那么，Poincaré 对偶映射
\[
HF^*(\widetilde\Delta_1^\#,\widetilde\Delta_0^\#)\cong HF^{n+1-*}(\widetilde\Delta_0^\#,\widetilde\Delta_1^\#)^\vee
\]
与 $\Z/2$ 作用反相容，即相差一个整体 $-1$ 符号而等变。
\end{lemma}
\begin{proof}
这些纯粹是上同调层次上的陈述，由\secref{14b}中列出的等变 Floer 上同调的一般性质推出。事实上，$\iota_{\widetilde E}|\widetilde\Delta$ 是 $\widetilde\Delta$ 的反定向自同构，而 PSS 同构与 $\Z/2$ 作用相容，故 (i) 成立。对于 (ii)，回顾 Poincaré 对偶同构来自三角形乘积
\[
\begin{aligned}
&HF^*(\widetilde\Delta_1^\#,\widetilde\Delta_0^\#)\otimes HF^{n+1-*}(\widetilde\Delta_0^\#,\widetilde\Delta_1^\#)\\
&\qquad\longrightarrow HF^{n+1}(\widetilde\Delta_0^\#,\widetilde\Delta_0^\#)\cong H^{n+1}(\widetilde\Delta_0;\K)\cong\K,
\end{aligned}
\]
因此，更典范的写法应为
\[
HF^*(\widetilde\Delta_1^\#,\widetilde\Delta_0^\#)\cong H^{n+1}(\widetilde\Delta_0;\K)\otimes HF^{n+1-*}(\widetilde\Delta_0^\#,\widetilde\Delta_1^\#)^\vee.
\]
在这一形式下，它仍与 $\Z/2$ 作用相容；由 (i)，去掉 $H^{n+1}(\widetilde\Delta_0;\K)$ 因子便引入一个符号。\znote{279}{两条要点。其一，(i) 中两种作用的差别就是「$\Id$ 与 $-\Id$ 不同」的全部内容：若 $\operatorname{char}\K=2$，引理退化为「作用平凡」，此后一切用 $\Z/2$-不变性区分次数、进而判定态射群无上同调（命题 18.14）的论证都会失效。其二，(ii) 的反相容有具体来源：Poincaré 对偶把 $HF^*$ 与 $HF^{n+1-*}$ 配对到 $HF^{n+1}(\widetilde\Delta_0^\#,\widetilde\Delta_0^\#)\cong\K$，而 $\Z/2$ 在这个基类上按 $-\Id$ 作用，故对偶映射整体带一个符号。二者与\eref{18.23}合用，便得 $j>k$ 时 $HF^*(\widetilde\Delta_j^\#,\widetilde\Delta_k^\#)$ 上的作用为 $-\Id$。[来源：原书 §18e、§18f（印刷页 279、282）]}
\end{proof}

如同 (U) 型膜的情形，特殊选择 Floer 数据和扰动数据，有助于使链层次上的 $\Z/2$ 作用更清楚。取两条消失路径 $\gamma_k$，相应的 (B) 型膜为 $\widetilde\Delta_k^\#$，消失圈为 $V_k^\#$。假设
\begin{equation}\label{eq:18.22}
\begin{cases}
\arg(\gamma'_0(0))>\arg(\gamma'_1(0)),\\
\gamma_0([0,1])\cap\gamma_1([0,1])=\{*\},\\
V_0\pitchfork V_1.
\end{cases}
\end{equation}
为 $(\widetilde\Delta_0,\widetilde\Delta_1)$ 选择 Floer 数据 $(H,J)$ 时，要施加比\stref{假设}{14.5}强得多的要求：对每个 $J_t$，$\widetilde\pi$ 都为 $J_t$-全纯，且 $H=0$。第二个条件与\stref{假设}{14.5}相容，因为由\eref{18.22}，法向角满足 $\sigma_1-\sigma_0\in(-\pi,0)$，故\eref{14.10}中可以取 $h=0$。任意交点 $\widetilde y=(y,x)\in\C(\widetilde\Delta_0,\widetilde\Delta_1)$ 在 $\widetilde\pi$ 下必投影到 $\widetilde\gamma_0$ 与 $\widetilde\gamma_1$ 的唯一交点。因此 $y=0$，意味着 $\widetilde y$ 位于 $\widetilde E^{\mathrm{fix}}=M$ 内。按定义，每个 $\widetilde\Delta_k$ 上的等变 Pin 结构，对应于 $\widetilde\Delta_k^{\mathrm{fix}}$ 的法丛平凡化，该平凡化投影到虚部为正的平方根 $\sqrt{\gamma'_k(0)}$。该平方根为 $\exp(i\sigma_k)=\exp(i\arg(\gamma'_k(0))/2)$ 的正实数倍。因此，\stref{引理}{14.11}中的符号 $w_k$ 都为 $+1$，意味着 $|o(\widetilde y)|_{\K}$ 上的 $\Z/2$ 作用平凡。这就说明，整个上链群 $CF^*(\widetilde\Delta_0^\#,\widetilde\Delta_1^\#)$ 都是 $\Z/2$-不变的。

类似地，任意 $\widetilde u\in\M_Z(\widetilde y_0,\widetilde y_1)$ 都投影到全纯映射 $\widetilde v=\widetilde\pi\circ\widetilde u:Z\to\widetilde S$，满足 $\widetilde v(\R\times\{k\})\subset\widetilde\gamma_k([-1,1])$、$\lim_{s\to\pm\infty}\widetilde v(s,\cdot)=0$。任何这样的 $\widetilde v$ 的能量都
\SourcePage{288}
必为零，因此它是常值映射，意味着 $\widetilde u$ 也完全落在 $M$ 中。假设 $(H,J)$ 的选择使 $M$ 上的 Floer 模空间正则，则\stref{引理}{14.6}蕴含 $\widetilde E$ 中的正则性。此外，\stref{引理}{14.9}说明，无论在 $M$ 还是在 $\widetilde E$ 中工作，绝对指标与定向空间都相同；\eref{14.13}对每个孤立的 $\widetilde u$ 贡献于 Floer 微分时所带的符号，也作出了相同结论。总之，存在复形同构
\begin{equation}\label{eq:18.23}
CF^*(\widetilde\Delta_0^\#,\widetilde\Delta_1^\#)\cong CF^*(V_0^\#,V_1^\#),
\end{equation}
其中右端消失圈上的膜结构是从 $\widetilde\Delta_k^\#$ 诱导的。此外，左端的典范 $\Z/2$ 作用平凡。更一般地，假设有消失路径 $\gamma_0,\ldots,\gamma_d$，$d\ge2$，满足
\begin{equation}\label{eq:18.24}
\begin{cases}
\arg(\gamma'_k(0))>\arg(\gamma'_{k+1}(0))\quad\text{对所有 }k,\\
\gamma_j([0,1])\cap\gamma_k([0,1])=\{*\}\quad\text{若 }j\ne k,\\
\text{各 $V_k$ 处于一般位置。}
\end{cases}
\end{equation}
那么，选取满足 $K=0$ 且 $\widetilde\pi$ 为 $J$-全纯的扰动数据 $(K,J)$ 后，就发现所有全纯 $(d+1)$ 边形都落在 $M$ 中。假设正则性成立，见\stref{引理}{14.7}，它们对 $M$ 与 $\widetilde E$ 上的 $\Ai$ 结构也给出相同贡献。

\StatementNumber{10}\begin{note}\label{st:18.10}
由\stref{引理}{18.9}(ii)，\eref{18.23}意味着，对偶群 $HF^*(\widetilde\Delta_1^\#,\widetilde\Delta_0^\#)$ 上的典范 $\Z/2$ 作用为 $-\Id$。原则上，可取非齐次项为零的 Floer 数据，在上链层次上证明这一点。为使其满足\eref{14.10}，其中必有 $h=0$，就须将法向角 $\sigma_1$ 换成 $\sigma_1+\pi$；这意味着\stref{引理}{14.11}中 $w_1=-1$，使相关定向空间上的 $\Z/2$ 作用非平凡。遗憾的是，法向角的改变使这一选择与此前对 $(\widetilde\Delta_0,\widetilde\Delta_1)$ 使用的 Floer 数据不相容。具体地，尝试同时使用这两者，会遇到伪全纯多边形的等变横截性方面的潜在问题，因此这里放弃这种做法。这意味着，在我们的 $\F(\widetilde E)^{\mathrm{naive}}$ 版本中，链层次上的等式\eref{18.23}成立，而 $(\widetilde\Delta_1,\widetilde\Delta_0)$ 的对应陈述只在上同调层次上成立。以下\stref{命题}{18.14}决定了这一选择。
\end{note}

\StatementNumber{11}\begin{note}\label{st:18.11}
在混合情形中，即一个子流形 $\widetilde L=\beta_{\widetilde E}^{-1}(L)$ 为 (U) 型，而另一个 $\widetilde\Delta=\beta_{\widetilde E}^{-1}(\Delta)$ 为 (B) 型时，类似于\eref{18.20}的结果成立。首先，虽然 $\Delta$ 有边界，仍可容易地定义 $(L,\Delta)$ 的 Floer 上同调，方法如下。取在 $E^{\mathrm{shrink}}$ 之外平凡的 Floer 数据 $(H,J)$，使任意 $u\in\M_Z(y_0,y_1)$ 对应的 $v=\pi\circ u$ 在 $v^{-1}(S^{\mathrm{shrink}})$ 之外全纯。此时，$v$ 的像不包含 $*$；否则，由开映射性，它也必须包含附近不属于 $\gamma([0,1])$ 的点 $z$，这意味着相交数 $u\cdot[E_z]$ 为
\SourcePage{289}
正。由于 $S\setminus(S^{\mathrm{shrink}}\cup\gamma([0,1]))$ 连通，同样的正性也会对 $z\in\partial S$ 成立，产生矛盾。结论是，由于 $u$ 不与 $\partial\Delta$ 相交，在 Floer 理论的考察中可以忽略该边界的存在。据此，可以利用同痕不变性，将问题归约到如下情形：$t\in[0,\delta]$ 时 $\gamma(t)=*-t$，而 $\gamma$ 的其余部分落在 $S^{\mathrm{shrink}}$ 内。此时，可如前对 $|v|$ 应用最大值原理，得到
\[
HF^*(\widetilde L^\#,\widetilde\Delta^\#)\cong HF^*(L^\#,\Delta^\#)\otimes_{\K}\K[+,-].
\]
在该等式右端，$\Delta$ 带有唯一的平凡 Pin 结构，分次则从 $\widetilde\Delta$ 的分次下降得到，即反过来应用\eref{18.18}的过程。
\end{note}
